\section*{अध्याय \ref{ch:b3:sequence-stratigraphy} --- अनुक्रम स्तरिकी}

\begin{solution}{exo:b3:sequence-stratigraphy:1}
$A = 0.5 + 0.1 = \qty{0.6}{m/kyr}$; \qty{20}{kyr} में \qty{12}{m} का
आवास-स्थान प्राप्त होता है। समुद्र-तल गिरने पर $A = -0.3 + 0.1 = \qty{-0.2}{m/kyr}$: एक
सापेक्ष पतन, जिसे अवधसान ने धीमा किया है।
\end{solution}

\begin{solution}{exo:b3:sequence-stratigraphy:2}
(क) \omterm{def:b3:sequence-stratigraphy:tract}{पतन-चरण तंत्र-पथ} (\omterm{def:b3:sequence-stratigraphy:regression}{बलात् रिग्रेशन})। (ख) \omterm{def:b3:sequence-stratigraphy:tract}{ट्रांसग्रेसिव तंत्र-पथ}। (ग)
\omterm{def:b3:sequence-stratigraphy:tract}{उच्च-स्तर तंत्र-पथ}। (घ) \omterm{def:b3:sequence-stratigraphy:tract}{निम्न-स्तर तंत्र-पथ} (पतन के बाद \omterm{def:b3:sequence-stratigraphy:regression}{सामान्य रिग्रेशन})।
\end{solution}

\begin{solution}{exo:b3:sequence-stratigraphy:3}
प्रत्येक \omterm{def:b3:sequence-stratigraphy:parasequence}{पैरा-अनुक्रम} अपने नीचे वाले से उथला समाप्त होता है: यह एक प्रगामी स्टैक है,
जिसमें तटरेखा समुद्र की ओर खिसकती है, जैसा उच्च-स्तर या \omterm{def:b3:sequence-stratigraphy:tract}{निम्न-स्तर तंत्र-पथ} में होता है। इसके
उलट क्रम, जिसमें प्रत्येक \omterm{def:b3:sequence-stratigraphy:parasequence}{पैरा-अनुक्रम} अधिक महीन और गहरा समाप्त होता है, प्रतिगामी है:
यह \omterm{def:b3:sequence-stratigraphy:tract}{ट्रांसग्रेसिव तंत्र-पथ} है।
\end{solution}

\begin{solution}{exo:b3:sequence-stratigraphy:4}
एक बंडल में पाँच युग्मक होना लघु उत्केंद्रता (लगभग
\qty{100}{kyr}) और पुरस्सरण (लगभग \qty{20}{kyr}) के अनुपात के बराबर है: प्रत्येक
युग्मक एक पुरस्सरण चक्र है, और प्रत्येक बंडल एक उत्केंद्रता चक्र। दर
$1.6/100 = \qty{1.6}{cm/kyr}$ है, और \qty{0.32}{m} मोटाई का एक युग्मक लगभग
\qty{20}{kyr} तक चलता है, जो सुसंगत है।
\end{solution}

\begin{solution}{exo:b3:sequence-stratigraphy:5}
$\mathrm dx_s/\mathrm dt = (4000 - 0.15\times20\,000)/50 = \qty{20}{m/kyr}$,
समुद्र की ओर। तटरेखा तब स्थिर रहती है जब $AW = Q$ हो: $A = 4000/20\,000 =
\qty{0.2}{m/kyr}$।
\end{solution}

\begin{solution}{exo:b3:sequence-stratigraphy:6}
$\mathrm dx_s/\mathrm dt = 4000/50 + 0.2/0.01 = 80 + 20 = \qty{100}{m/kyr}$;
केवल पतन \qty{20}{m/kyr} का हिसाब रखता है, यानी पाँचवाँ भाग।
\end{solution}

\begin{solution}{exo:b3:sequence-stratigraphy:7}
यह तब पीछे हटती है जब $Q < AW = 0.4\times25\,000 = \qty{10000}{m^3/kyr}$ प्रति
मीटर। $Q = 6000$ के लिए: $(6000 - 10\,000)/10 = \qty{-400}{m/kyr}$, यानी
हर हज़ार वर्ष में \qty{0.4}{km} का पीछे हटना होता है: निम्न तटमुख
तीव्र ट्रांसग्रेशन बनाता है।
\end{solution}

\begin{solution}{exo:b3:sequence-stratigraphy:8}
इन्हें अलग करने के लिए $1/94.9 - 1/123.9 = \num{0.00247}\text{ चक्र}/\unit{kyr}$:
$1/0.00247 = \qty{406}{kyr}$ से लंबा अभिलेख चाहिए, और \qty{600}{kyr} (विभेदन
0.00167) पर्याप्त है। $1/94.9 - 1/98.9 = 0.00042$: इसके लिए लगभग
\qty{2400}{kyr} चाहिए, जो अभिलेख से चार गुना लंबा है।
\end{solution}

\begin{solution}{exo:b3:sequence-stratigraphy:9}
$16.4/(20\times40.98) = \qty{0.0200}{m/kyr}$, \qty{2.0}{cm/kyr}; राख संस्तर
$3.1/0.020 = \qty{155}{kyr}$ की दूरी पर हैं।
\end{solution}

\begin{solution}{exo:b3:sequence-stratigraphy:10}
(क) $\tan\alpha = -0.4/200 = -0.002$: \ang{0.11} पर नीचे की ओर झुकता है। (ख)
$0.2/50 = 0.004$: \ang{0.23} पर ऊपर की ओर चढ़ता है। भूकंपीय खंड दसियों
किलोमीटर फैलता है, ऊर्ध्वाधर अतिरंजन दस से सौ गुना है, इसलिए कोण दिखता है;
कगार, जो कुछ सौ मीटर लंबा है, एक मीटर या उससे कम का उत्थान दिखाता है, जो
संस्तरों के बदलाव में खो जाता है।
\end{solution}

\begin{solution}{exo:b3:sequence-stratigraphy:11}
$1\,296\,000/54.725282 = \qty{23.68}{kyr}$ ($g_5$); $/57.924383 =
\qty{22.37}{kyr}$ ($g_2$); $/68.377912 = \qty{18.95}{kyr}$ ($g_4$);
$/67.834313 = \qty{19.11}{kyr}$ ($g_3$). $g_2 - g_5 = 3.199101''$ प्रति वर्ष,
\qty{405.1}{kyr}; $g_4 - g_5 = 13.65263''$, \qty{94.9}{kyr}।
\end{solution}

\begin{solution}{exo:b3:sequence-stratigraphy:12}
दोनों एक ही दर पर $0.3\times100 = \qty{30}{m}$ आवास-स्थान बनाते हैं; समान
आपूर्ति के साथ ज्यामिति एक जैसी रहती है, क्योंकि अवसाद केवल सापेक्ष समुद्र-तल
पर प्रतिक्रिया करता है। यूस्टेटिक परिवर्तन हर जगह एक समान होता है: दूसरे
किनारों पर उसी आयु का वही अनुक्रम पहले इतिहास की ओर संकेत करता है। बेसिन
के ऊष्मीय और भार इतिहास से बना अवधसान इतिहास
(बैकस्ट्रिपिंग), या ऑक्सीजन समस्थानिक जैसा कोई स्वतंत्र यूस्टेटिक अभिलेख,
इन्हें भी अलग करता है।
\end{solution}

\begin{solution}{pb:b3:sequence-stratigraphy:1}
\textbf{1.} $\dot E = 30\times2\pi/400\times\sin(2\pi t/400)$: अधिकतम
\qty{0.471}{m/kyr}, \qty{100}{kyr} और \qty{500}{kyr} पर बढ़ते हुए, और
\qty{300}{kyr} पर घटते हुए।

\textbf{2.} वलन-मोड़ पर \qty{0.020}{m/kyr}, \qty{20}{km} पर 0.037, और
\qty{50}{km} पर 0.062, \qty{120}{km} पर \qty{0.120}{m/kyr}।

\textbf{3.} $A < 0$ तब होता है जब $\sin(2\pi t/400) < -0.0346/0.471 = -0.0734$:\newline $t =
200 + (400/2\pi)\arcsin0.0734 = \qty{204.7}{kyr}$।

\textbf{4.} $\sigma = 0.0617$ के लिए: $t = 400 - 63.66\arcsin0.131 =
\qty{391.6}{kyr}$।

\textbf{5.} लगभग \qty{187}{kyr} बनाम 200: यूस्टेटिक पतन को पहले अवधसान को
पार करना होगा, और निम्न-स्तर से पहले वह उससे आगे निकलना बंद कर देता है।

\textbf{6.} $0.045\times400 = \qty{18}{m}$, यूस्टेटिक परास \qty{60}{m} के
मुकाबले: एक तिहाई।

\textbf{7.} 50 से लगभग \qty{120}{kyr} तक भूमि की ओर, लगभग
\qty{440}{kyr} तक समुद्र की ओर, लगभग \qty{560}{kyr} तक भूमि की ओर, फिर
\qty{600}{kyr} तक फिर समुद्र की ओर (मॉडल के पलटाव बिंदु 118.5, 439.5 और
\qty{559}{kyr} हैं; हर \qty{50}{kyr} पर बनी सारणी उन्हें केवल सीमांकित करती है)।

\textbf{8.} $(6000 - 0.30\times15\,000)/40 = \qty{37.5}{m/kyr}$, समुद्र की ओर।

\textbf{9.} $6000/15\,000 = \qty{0.40}{m/kyr}$.

\textbf{10.} $6000/35.7 + 0.424/0.012 = 168.1 + 35.3 = \qty{203}{m/kyr}$,
मॉडल से \qty{0.5}{\percent} के भीतर।

\textbf{11.} आपूर्ति: गति का \qty{83}{\percent}।

\textbf{12.} $(-42.7 - 16.3)/(50.7 - 17.4) = \qty{-1.77}{m/km}$, \ang{0.10}।
तब से पहले बिंदु का अवधसान हुआ है: $0.0345\times395.5 =
\qty{13.6}{m}$, दूसरे का $0.0623\times208.5 = \qty{13.0}{m}$: लगभग
समान है, इसलिए वर्तमान ऊँचाइयों में \qty{59.0}{m} की गिरावट
समुद्र-तल के पतन से केवल \qty{0.6}{m} भिन्न है। उस समय तटरेखा पर सापेक्ष
पतन, यानी पतन के दौरान समुद्र-तल के पतन में से अवधसान घटाने पर, कम था।

\textbf{13.} 150 और 200: उच्च-स्तर; 250, 300, 350: पतन-चरण; 400:
निम्न-स्तर; 450, 500, 550: ट्रांसग्रेसिव (तटरेखा 440 के ठीक बाद
पीछे हटती है); 600: उच्च-स्तर।

\textbf{14.} 204.7 और \qty{391.6}{kyr} बनाम 204.5 और 391.5: समान,
मॉडल के \qty{0.5}{kyr} चरण के भीतर।

\textbf{15.} शेल्फ़ पर, पतन के दौरान उजागर सतह: उच्च-स्तर निक्षेपों का शीर्ष
वलन-मोड़ से \qty{17}{km} तक फैला है और अधोपातन करते पतन-चरण फ़ोरसेट के
ऊपर से आगे जाता है; द्रोणी में सहसंबंधी
संगति, परिपाटी के अनुसार, \omterm{def:b3:sequence-stratigraphy:tract}{पतन-चरण तंत्र-पथ} के आधार या शीर्ष पर होती है। वास्तविक
शेल्फ़ नदियों द्वारा कटकर कटी हुई घाटियों में बदल जाता, मिट्टी में अपक्षयित होता,
और आंशिक रूप से अपरदित होता।

\textbf{16.} $(18.5 - 52.8)/(559 - 439.5) = \qty{-0.29}{km/kyr}$; \omterm{def:b3:sequence-stratigraphy:surfaces}{अधिकतम
प्लावन सतह} उसके अंत पर बनती है।

\textbf{17.} तटरेखा लगभग समतल तटीय मैदान पर तेज़ी से खिसकती है:
आपूर्ति उसके पास फँस जाती है, विस्तृत क्षेत्र में फैलती है, और हर स्थान
को थोड़ी ही प्राप्त होती है। पाद के आगे, लगभग \qty{55}{km} से परे, समुद्र-तल
अभावग्रस्त है: वहाँ \omterm{def:b3:sequence-stratigraphy:condensed}{संघनित काट} बनती है।

\textbf{18.} विभेदन $1/30 = \num{0.033}\text{ चक्र}/\unit{m}$; नाइक्विस्ट $1/(2\times
0.01) = \num{50}\text{ चक्र}/\unit{m}$।

\textbf{19.} $7.5/2.0 = 3.75$, 1 और $0.81/2.0 = 0.405$; कक्षीय अनुपात
$405/100 = 4.05$ और $41/100 = 0.41$ हैं।

\textbf{20.} \qty{100}{kyr} में \qty{2.0}{m}: \qty{2.0}{cm/kyr}; $0.81/41 =
\qty{1.98}{cm/kyr}$; $7.5/405 = \qty{1.85}{cm/kyr}$। \qty{8}{m} के निकट
आवर्तलेख जो तरंगदैर्घ्य दिखा सकता है, वे 10, 7.5 और \qty{6}{m} हैं: सच्चा
\qty{8.1}{m} उनमें से दो के बीच आता है।

\textbf{21.} $30/0.020 = \qty{1500}{kyr}$।

\textbf{22.} $23.7\times0.020 = \qty{0.474}{m}$। हर
\qty{25}{cm} पर नमूना लेने पर नाइक्विस्ट आवृत्ति प्रति मीटर 2 चक्र है, तरंगदैर्घ्य
\qty{0.5}{m}: पुरस्सरण, 2.11 चक्र प्रति मीटर पर, उससे परे है और
$4 - 2.11 = 1.89$ पर वापस मुड़ता है, यानी \qty{0.53}{m} का मिथ्या तरंगदैर्घ्य।

\textbf{23.} कोई भी एकल शिखर, चाहे उसका मूल कुछ भी हो, किसी दर पर किसी आवर्तकाल से
मेल खा जाता है; तीन शिखर जिनके अनुपात कक्षीय अनुपातों से मेल खाते हैं, दर से स्वतंत्र हैं,
और तब दर हर एक से समान निकलनी चाहिए।

\textbf{24.} तटरेखा पतन के आरंभ में \qty{17.37}{km} से
निम्न-स्तर रिग्रेशन के अंत में \qty{52.84}{km} तक खिसकती है: अधिकतम
द्रोणी-अभिमुख विस्थापन \textbf{\qty[text-series-to-math=true]{35.5}{km}}।
\end{solution}
